% !TeX program = xelatex
% Windows 编译准备：py -m pip install Pygments
% 编译命令（运行两次）：xelatex -shell-escape report.tex

\documentclass[UTF8,12pt,a4paper,fontset=windows]{ctexart}

\usepackage{geometry}
\geometry{left=2.6cm,right=2.6cm,top=2.7cm,bottom=2.7cm}
\usepackage[table]{xcolor}
\usepackage{graphicx}
\usepackage{booktabs}
\usepackage{tabularx}
\usepackage{array}
\usepackage{amsmath}
\usepackage{enumitem}
\usepackage{float}
\usepackage{caption}
\usepackage{fancyhdr}
\usepackage{minted}
\usepackage[hidelinks]{hyperref}

\definecolor{PrimaryBlue}{HTML}{1F4E79}
\definecolor{CodeBackground}{HTML}{F7F8FA}
\definecolor{CodeFrame}{HTML}{AEB8C2}
\definecolor{TableHeader}{HTML}{DCE6F1}
\definecolor{MutedText}{HTML}{5B6570}

\newcolumntype{Y}{>{\raggedright\arraybackslash}X}
\newcolumntype{C}{>{\centering\arraybackslash}X}
\renewcommand{\arraystretch}{1.28}

\setminted{
  linenos,
  frame=single,
  framesep=2mm,
  rulecolor=\color{CodeFrame},
  bgcolor=CodeBackground,
  fontsize=\small,
  breaklines,
  breakanywhere,
  autogobble,
  tabsize=4,
  numbersep=8pt,
  xleftmargin=1.4em
}
\renewcommand{\theFancyVerbLine}{%
  \sffamily\scriptsize\color{MutedText}\arabic{FancyVerbLine}}

\ctexset{
  section={
    name={,、},
    number=\chinese{section},
    format=\Large\bfseries\color{PrimaryBlue}
  },
  subsection={
    format=\large\bfseries
  },
  subsubsection={
    format=\normalsize\bfseries
  }
}

\pagestyle{fancy}
\fancyhf{}
\fancyhead[L]{计算机网络实验}
\fancyhead[R]{互联网体系结构初识}
\fancyfoot[C]{\thepage}
\setlength{\headheight}{15pt}

\captionsetup{font=small,labelfont=bf}
\setlist{itemsep=0.35em,topsep=0.45em}
\setlength{\parindent}{2em}
\setlength{\parskip}{0.25em}

\begin{document}

\begin{center}
  {\zihao{2}\bfseries\color{PrimaryBlue} 互联网体系结构初识实验报告\par}
  \vspace{0.8em}
  \begin{tabularx}{\textwidth}{CCCC}
    姓名：吴峥 & 学号：2024K8009909013 & 班级：2406 & 日期：2026 年 9 月 4 日
  \end{tabularx}
\end{center}
\vspace{0.3em}
\hrule
\vspace{1.2em}

\section{实验目的}

本实验通过 Wireshark 观察访问网站时产生的数据包，初步了解 ARP、DNS、TCP、HTTP 和 HTTPS 等协议的作用，以及协议之间的分层与封装关系。同时，使用 Python 对多个网站进行访问测量，将实际访问时间与由地理距离估算的理论传播时延进行比较，认识真实互联网时延的主要来源。

\section{实验环境}

本地实验在 Windows Subsystem for Linux 2（WSL 2）中完成，具体环境如表~\ref{tab:environment} 所示。

\begin{table}[H]
  \centering
  \caption{实验环境}
  \label{tab:environment}
  \begin{tabularx}{\textwidth}{>{\bfseries}p{3.2cm}Y}
    \toprule
    \rowcolor{TableHeader} 项目 & 版本或说明 \\
    \midrule
    WSL & WSL 2，程序版本 2.7.11.0 \\
    Linux 发行版 & Ubuntu 24.04.4 LTS（Noble Numbat） \\
    Linux 内核 & \texttt{6.18.33.2-microsoft-standard-WSL2}，x86-64 \\
    Mininet & 2.3.0 \\
    Wireshark & 4.2.2 \\
    Python & 3.14.6 \\
    主要 Python 库 & \texttt{pycurl 7.47.0}、\texttt{geoip2 5.3.0}、\texttt{matplotlib 3.11.1} \\
    其他工具与数据 & \texttt{wget}、GeoLite2 City IP 地理位置数据库 \\
    \bottomrule
  \end{tabularx}
\end{table}

其中，Mininet 用于建立能够访问互联网的虚拟主机环境，Wireshark 用于抓取和观察数据包，Python 脚本负责网站性能测量、数据清洗和绘图。较大规模的网站测量后来使用相同脚本在阿里云远程 Linux 主机上完成。

\section{互联网协议观察实验}

\subsection{实验步骤}

首先启动能够访问互联网的 Mininet 环境：

\begin{minted}{bash}
sudo mn --nat
\end{minted}

在 Mininet 中打开主机 \texttt{h1} 的终端：

\begin{minted}{text}
mininet> xterm h1
\end{minted}

在 \texttt{h1} 中配置 DNS 服务器并启动 Wireshark：

\begin{minted}{bash}
echo "nameserver 8.8.8.8" > /etc/resolv.conf
wireshark &
\end{minted}

开始抓包后，使用 \texttt{wget} 下载中国科学院大学主页：

\begin{minted}{bash}
wget -U "Mozilla/5.0 (Windows NT 10.0; Win64; x64) AppleWebKit/537.36 (KHTML, like Gecko) Chrome/123.0 Safari/537.36" https://www.ucas.ac.cn
\end{minted}

停止抓包并保存为 \texttt{download-wire-shark.pcapng}。该文件包含 75 个数据包，抓包时长约为 5.14 秒。

\subsection{各协议的简单说明}

\begin{table}[H]
  \centering
  \small
  \caption{实验涉及的主要协议}
  \begin{tabularx}{\textwidth}{>{\bfseries}p{1.7cm}p{5.0cm}Y}
    \toprule
    \rowcolor{TableHeader} 协议 & 主要作用 & 本次实验中的表现 \\
    \midrule
    ARP & 根据同一链路中的 IP 地址查询对应的 MAC 地址 & 出现了 ``who has'' 请求和 ``is at'' 应答 \\
    DNS & 将容易记忆的域名解析为 IP 地址 & 将 \texttt{www.ucas.ac.cn} 解析为 \texttt{124.16.77.5} \\
    TCP & 在通信双方之间建立可靠的字节流连接 & 与服务器的 443 端口完成握手、传输和关闭连接 \\
    HTTP & 规定下载工具与 Web 服务器之间的请求和响应格式 & 本次请求使用 HTTPS，因此 HTTP 内容没有以明文显示 \\
    HTTPS & 在 HTTP 与 TCP 之间加入 TLS 加密和身份验证 & Wireshark 中可以看到 TLS 握手及后续加密数据 \\
    \bottomrule
  \end{tabularx}
\end{table}

\subsection{协议的分层与封装}

抓包中可以观察到不同协议按层次封装。例如：

\begin{minted}{text}
Ethernet < IP < UDP < DNS
Ethernet < IP < TCP < TLS < HTTPS 数据
\end{minted}

发送数据时，上层内容逐层加上首部后交给下层；接收数据时则按相反顺序逐层解析。DNS 查询通常使用 UDP，而 HTTPS 通过 TCP 可靠传输。这里的 HTTPS 实际上是经过 TLS 加密的 HTTP，因此抓包能够识别 TLS 记录，却不能直接看到加密后的 HTTP 请求正文。

\subsection{抓包结果}

抓包中的主要通信过程如表~\ref{tab:capture-overview} 所示。

\begin{table}[H]
  \centering
  \small
  \caption{抓包通信过程概览}
  \label{tab:capture-overview}
  \begin{tabularx}{\textwidth}{p{2.8cm}p{2.2cm}Y}
    \toprule
    \rowcolor{TableHeader} 阶段 & 数据包编号 & 观察结果 \\
    \midrule
    DNS 查询 & 1--4 & \texttt{h1} 同时查询域名的 IPv4（A）和 IPv6（AAAA）地址，A 记录返回 \texttt{124.16.77.5} \\
    TCP 建立连接 & 5--7 & 客户端发送 SYN，服务器返回 SYN/ACK，客户端再发送 ACK，完成三次握手 \\
    TLS 握手 & 8--18 & 双方协商加密参数并交换服务器证书，为 HTTPS 通信建立安全通道 \\
    网页数据传输 & 19--66 & 服务器向客户端连续发送加密的网页数据，客户端发送 ACK 确认收到的数据 \\
    TCP 关闭连接 & 67--69 & 双方使用 FIN/ACK 结束本次连接 \\
    ARP 交互 & 70--75 & 网络节点查询并应答局域网 IP 地址对应的 MAC 地址 \\
    \bottomrule
  \end{tabularx}
\end{table}

DNS 查询从 \texttt{10.0.0.1} 发往 DNS 转发节点 \texttt{10.255.255.254}，约 30 ms 后收到了 IPv4 地址响应。随后 \texttt{10.0.0.1} 使用临时端口 \texttt{53446} 连接服务器 \texttt{124.16.77.5} 的 HTTPS 端口 \texttt{443}。TCP 三次握手完成后进行 TLS 握手，之后 \texttt{wget} 的 HTTP 请求和服务器返回的页面内容都以加密数据形式传输。

ARP 数据包出现在抓包末尾，说明开始下载时所需的 MAC 地址很可能已经存在于缓存中，因而不需要在 TCP 连接之前重新查询。后续 ARP 请求仍然展示了 IP 地址到 MAC 地址的解析过程。

\subsubsection{TCP 三次握手}

建立 TCP 连接的过程称为“三次握手”。抓包开头出现的几个方向交替的数据包，含义如下。设客户端的初始序号为 $x$，服务器的初始序号为 $y$：

\begin{table}[H]
  \centering
  \small
  \caption{TCP 三次握手}
  \begin{tabularx}{\textwidth}{C p{3.0cm}p{2.1cm}p{3.1cm}Y}
    \toprule
    \rowcolor{TableHeader} 编号 & 方向 & 标志位 & 序号与确认号 & 含义 \\
    \midrule
    5 & 客户端 $\rightarrow$ 服务器 & SYN & \texttt{Seq=x} & 客户端请求建立连接，并给出自己的初始序号 \\
    6 & 服务器 $\rightarrow$ 客户端 & SYN, ACK & \texttt{Seq=y, Ack=x+1} & 服务器同意连接，确认客户端的请求，同时给出自己的初始序号 \\
    7 & 客户端 $\rightarrow$ 服务器 & ACK & \texttt{Seq=x+1, Ack=y+1} & 客户端确认服务器的请求，连接建立完成 \\
    \bottomrule
  \end{tabularx}
\end{table}

从第 5 个包发出 SYN 到第 7 个包完成确认大约经过 2.36 ms。SYN 本身会占用一个序号，因此双方确认时都使用“对方初始序号加 1”。握手数据包中还协商了最大报文段长度（MSS）、选择确认（SACK）和窗口扩大等参数，本次只需要知道这些参数用于提高后续传输效率即可。

\subsubsection{TLS 握手}

TCP 连接只能保证数据可靠到达，并不能保证内容不被窃听或修改，因此访问 HTTPS 网站还需要进行 TLS 握手。

\begin{table}[H]
  \centering
  \small
  \caption{TLS 握手过程}
  \begin{tabularx}{\textwidth}{C p{4.2cm}Y}
    \toprule
    \rowcolor{TableHeader} 编号 & 主要内容 & 简单说明 \\
    \midrule
    8 & Client Hello & 客户端给出支持的 TLS 参数和加密算法 \\
    9 & TCP ACK & 服务器先确认已经收到 Client Hello \\
    10--14 & Server Hello、证书等 & 服务器选择通信参数并向客户端发送身份证书；内容较长，因此被拆成多个 TCP 段 \\
    15 & TCP ACK & 客户端确认已收到服务器发送的握手数据 \\
    16 & 客户端密钥交换与完成消息 & 客户端生成后续加密通信所需的信息 \\
    17 & 服务器完成消息 & 服务器确认 TLS 安全通道已经建立 \\
    18 & 加密的应用数据 & 客户端开始发送经过加密的 HTTP 请求 \\
    \bottomrule
  \end{tabularx}
\end{table}

这也解释了为什么 Wireshark 没有直接显示 \texttt{GET} 请求和网页正文：它们位于 TLS 加密数据内部。Wireshark 仍然可以显示通信方向、数据长度和 TLS 记录类型，但在没有会话密钥的情况下不能还原其中的 HTTP 明文。

\subsubsection{数据传输时的交替数据包}

第 19--66 个包主要是服务器发送网页数据、客户端返回 ACK。服务器并不一定每发送一个数据段就等待一次确认，它可以连续发送多个段；客户端的 ACK 也可以一次累计确认此前收到的所有连续数据。因此列表中会看到多个服务器数据包与客户端 ACK 交替出现，这实际上是服务器持续发送加密网页、客户端持续确认的过程。

\subsubsection{TCP 连接关闭}

结束连接的过程通常称为“四次挥手”。本次抓包只显示了 3 个包，是因为服务器把中间的“确认客户端 FIN”和“发送自己的 FIN”合并在同一个 \texttt{FIN, ACK} 包中。

\begin{table}[H]
  \centering
  \small
  \caption{TCP 连接关闭过程}
  \begin{tabularx}{\textwidth}{C p{3.2cm}p{2.2cm}Y}
    \toprule
    \rowcolor{TableHeader} 编号 & 方向 & 标志位 & 含义 \\
    \midrule
    67 & 客户端 $\rightarrow$ 服务器 & FIN, ACK & 客户端表示网页已经接收完成，不再发送新的应用数据 \\
    68 & 服务器 $\rightarrow$ 客户端 & FIN, ACK & 服务器确认客户端的 FIN，同时表示自己也不再发送新的应用数据 \\
    69 & 客户端 $\rightarrow$ 服务器 & ACK & 客户端确认服务器的 FIN，连接关闭 \\
    \bottomrule
  \end{tabularx}
\end{table}

因此，从逻辑上仍然包含双方各自的 FIN 和 ACK，只是其中两个动作被合并了。这种“三个数据包完成四个逻辑动作”的情况是正常的 TCP 连接关闭方式。

\subsection{网页下载过程总结}

从应用程序的角度看，下载网页大致经历了以下步骤：

\begin{enumerate}
  \item \texttt{wget} 根据 URL 得到域名 \texttt{www.ucas.ac.cn}。
  \item DNS 将域名解析为服务器 IP 地址 \texttt{124.16.77.5}。
  \item 主机利用 ARP 获得下一跳设备的 MAC 地址，在以太网中发送数据帧。
  \item 客户端与服务器的 443 端口进行 TCP 三次握手。
  \item 双方进行 TLS 握手，建立加密连接并验证服务器身份。
  \item 客户端发送加密的 HTTP 请求，服务器返回加密的网页内容。
  \item 下载完成后，双方关闭 TCP 连接。
\end{enumerate}

这个过程说明，一次看似简单的网页访问实际上需要应用层、传输层、网络层和链路层的多个协议共同工作。

\section{传输性能测量实验}

\subsection{测量方法}

实验脚本读取网站列表，对每个网站执行访问和延迟测量，记录以下四个指标：

\begin{itemize}
  \item \textbf{Ping：}向实际连接的服务器 IP 发送 3 次 ICMP Echo，取最小往返时间。
  \item \textbf{DNS：}完成域名解析所需的时间。
  \item \textbf{TCP transfer：}从收到 HTTP 响应的第一个字节到接收完最后一个字节的时间。
  \item \textbf{Total time：}从开始请求到完整接收响应的总时间。
\end{itemize}

脚本还使用 GeoIP 数据库估算服务器的地理位置，并计算测量点到服务器的大圆距离。将信号按真空光速往返传播所需的时间记为 \texttt{c-latency}：

\begin{equation}
  c\text{-latency}=\frac{2\times \text{distance}}{c},
\end{equation}

其中 $c=299792.458\ \mathrm{km/s}$。\texttt{c-latency} 是只考虑地理距离的理想下界，实际网络时延通常会大于它。

\subsection{测量脚本设计}

脚本按照“读取站点---访问测量---查询地理位置---保存结果---清洗绘图”的顺序工作：

\begin{minted}{text}
sites.csv
    ↓
读取并规范化域名
    ↓
并发执行 HTTP/HTTPS 访问和 Ping
    ↓
用 GeoIP 查询服务器位置并计算 c-latency
    ↓
逐条写入 CSV
    ↓
清洗无效记录并绘制 CDF
\end{minted}

访问测量主要使用 \texttt{pycurl} 提供的计时值。核心计算可简化为：

\begin{minted}{python}
dns_s = curl.getinfo(pycurl.NAMELOOKUP_TIME)
first_byte_s = curl.getinfo(pycurl.STARTTRANSFER_TIME)
total_s = curl.getinfo(pycurl.TOTAL_TIME)

dns_ms = dns_s * 1000
tcp_transfer_ms = (total_s - first_byte_s) * 1000
total_time_ms = total_s * 1000
\end{minted}

这里将“第一个响应字节到达以后，直到响应接收完成”的时间作为 TCP data transfer。Ping 部分使用系统的 \texttt{ping} 命令向网站实际连接的 IP 发送 3 次请求，并读取最小 RTT。这样测得的 Ping 与 HTTP 请求访问的是同一个服务器地址，便于比较。

地理距离使用两点经纬度计算，之后得到理想传播下界：

\begin{minted}{python}
distance_km = haversine_km(source_lat, source_lon,
                           server_lat, server_lon)
c_latency_ms = 2 * distance_km / 299_792.458 * 1000
\end{minted}

由于要测量大量网站，脚本使用线程池并发执行任务，默认并发数为 20。每完成一个网站就立即将结果写入 CSV 并刷新文件；如果测量意外中断，再次运行时可以跳过已经完成的域名。这一设计既缩短了总运行时间，也避免长时间测量因一次中断而全部丢失。

绘图前还要对结果进行基本清洗。程序检查 IP 地址、HTTP 状态、各项时间及 GeoIP 定位信息，最后分别计算四个指标与 c-latency 的比值，排序后绘制经验累计分布 CDF。具体清洗规则在下一节说明。

\subsection{实验过程与数据清洗}

\subsubsection{测量执行}

脚本使用北京附近的坐标作为测量点位置；最终的大规模测量在阿里云远程 Linux 主机上执行，同一脚本也可以在 Mininet 主机中运行。等价的主要命令如下：

\begin{minted}{bash}
python measure_websites.py \
  --sites sites.csv \
  --geoip-db GeoLite2-City.mmdb \
  --source-lat 39.9075 \
  --source-lon 116.3972 \
  --limit 20000 \
  --output results-20000.csv
\end{minted}

实际上进行了分阶段实验：本地先测试 1000 个网站，严格清洗数据时保留的数据在 100 多条，于是尝试 5000 条，大约半小时才完成，保留数据 700 多条。一方面绘制出的图形相对阶梯状明显，而且出现很明显的波浪式上升，猜测是数据量不够。最终是开了一个托管在阿里云的远程机跑了三个多小时，得到接近 3000 条有效数据，图形就相对平滑了。

\subsubsection{数据清洗逻辑}

原始数据包含 20,000 个网站。由于部分网站无法访问、拒绝 Ping、返回错误状态码，或者缺少可靠的 GeoIP 信息，不能直接把所有记录用于比较。清洗脚本按以下顺序检查每条数据：

\begin{enumerate}
  \item 五个数值字段 c-latency、Ping、DNS、TCP transfer 和 Total time 必须完整。
  \item 目标地址必须是有效的公网 IP。
  \item HTTP 状态码必须在 200--399 范围内，只保留正常响应或重定向。
  \item 各项数值必须大于 0，且 Total time 不能小于 DNS 与 TCP transfer 之和，否则各阶段时间关系不合理。
  \item Ping 不应小于按真空光速计算的 c-latency；如果出现这种情况，通常说明 GeoIP 位置不准确。
  \item GeoIP 给出的定位精度半径不能超过 200 km，并且测量点到服务器的距离应大于该精度半径。
\end{enumerate}

本次清洗结果如表~\ref{tab:cleaning} 所示。每条被剔除的记录只统计其首先发现的主要原因。

\begin{table}[H]
  \centering
  \small
  \caption{数据清洗结果}
  \label{tab:cleaning}
  \begin{tabularx}{\textwidth}{p{4.0cm}C Y}
    \toprule
    \rowcolor{TableHeader} 主要原因 & 剔除数量 & 说明 \\
    \midrule
    关键数值缺失 & 12,351 & 网站访问失败、Ping 无响应或缺少可用的地理位置等 \\
    GeoIP 定位过于粗略 & 2,856 & 定位精度半径超过 200 km，不适合计算较精确的地理下界 \\
    HTTP 状态不在 2xx/3xx & 1,671 & 网站返回拒绝、错误或其他非成功状态 \\
    Ping 小于 c-latency & 169 & 实测值违反该地理位置对应的理论传播下界 \\
    数值为 0 或负数 & 44 & 计时值不能用于倍数计算 \\
    \midrule
    \textbf{合计剔除} & \textbf{17,091} & --- \\
    \textbf{最终保留} & \textbf{2,909} & 五项指标完整、状态正常且地理信息较可靠 \\
    \bottomrule
  \end{tabularx}
\end{table}

\subsection{测量结果}

图~\ref{fig:cdf} 的横轴表示“实测时延除以 c-latency”，即实际网络相对于理想传播下界慢了多少倍；纵轴是累计分布 CDF。例如，CDF 等于 0.5 处的横坐标就是该指标的中位数。

\begin{figure}[H]
  \centering
  \includegraphics[width=\textwidth]{ping_latency_graph/latency_inflation_cdf_20000_clean_geo200_distance_ecdf.png}
  \caption{网站时延相对 c-latency 的累计分布}
  \label{fig:cdf}
\end{figure}

清洗后 2,909 个样本的主要结果如表~\ref{tab:medians} 所示。

\begin{table}[H]
  \centering
  \small
  \caption{各项指标的中位数}
  \label{tab:medians}
  \begin{tabularx}{\textwidth}{p{3.4cm}p{3.3cm}p{3.6cm}Y}
    \toprule
    \rowcolor{TableHeader} 指标 & 实测时间中位数 & 相对 c-latency 的中位倍数 & 简单理解 \\
    \midrule
    Ping RTT & 185.729 ms & 4.0$\times$ & 最接近单纯的网络往返传播时间 \\
    DNS lookup & 435.765 ms & 11.1$\times$ & 还受到 DNS 服务器和缓存状态影响 \\
    TCP data transfer & 159.109 ms & 4.5$\times$ & 受到响应大小、带宽和拥塞等影响 \\
    Total HTTP time & 2401.232 ms & 58.6$\times$ & 包含解析、建连、TLS、服务器处理和数据下载等完整过程 \\
    \bottomrule
  \end{tabularx}
\end{table}

为了进一步观察分布而不仅仅比较中位数，可以读取曲线在不同累计比例处的横坐标，结果见表~\ref{tab:percentiles}。

\begin{table}[H]
  \centering
  \small
  \caption{各项指标的倍数分位数}
  \label{tab:percentiles}
  \begin{tabularx}{\textwidth}{p{3.4cm}CCCC}
    \toprule
    \rowcolor{TableHeader} 指标 & 25\% & 50\% & 75\% & 90\% \\
    \midrule
    Ping RTT & 3.0$\times$ & 4.0$\times$ & 5.2$\times$ & 8.0$\times$ \\
    DNS lookup & 4.8$\times$ & 11.1$\times$ & 26.1$\times$ & 58.9$\times$ \\
    TCP data transfer & 0.03$\times$ & 4.5$\times$ & 13.1$\times$ & 33.1$\times$ \\
    Total HTTP time & 32.3$\times$ & 58.6$\times$ & 125.8$\times$ & 219.5$\times$ \\
    \bottomrule
  \end{tabularx}
\end{table}

\subsubsection{图中趋势和分布规律}

横轴采用对数坐标，因此曲线越靠左，表示该指标越接近理想传播下界；曲线上升越陡，表示多数样本越集中；向右延伸较长，则表示不同网站之间的差异较大。

\begin{itemize}
  \item \textbf{Ping RTT（蓝色）最靠左且上升较陡。}它的中位数约为 4.0$\times$，90\% 的样本不超过 8.0$\times$。这说明 Ping 的分布相对集中，主要反映网络路径上的传播、转发和排队时延。
  \item \textbf{DNS lookup（黑色）比 Ping 更靠右，也更分散。}中位数为 11.1$\times$，到 90\% 分位时达到 58.9$\times$。不同域名的缓存状态、DNS 服务器位置以及是否需要继续向上游查询都会造成较大的差别。
  \item \textbf{TCP data transfer（灰色）在左端有明显的一段快速上升，之后又形成较长的右尾。}约 34.7\% 的样本小于 1$\times$，绘图时被集中显示在 1$\times$ 的左边界。出现这种现象并不表示数据传播超过光速，因为该指标只计算“第一个响应字节到最后一个响应字节”的间隔；对于响应很小的网站，这段时间可以远小于一次完整 RTT。另一方面，大网页、低带宽或拥塞会把部分样本推向右侧，因此它的离散程度较大。
  \item \textbf{Total HTTP time（橙色）整体最靠右。}其中位数达到 58.6$\times$，90\% 的样本不超过约 219.5$\times$，并且存在明显长尾。它包含一次网页访问的多个阶段，任何一个阶段变慢都会累积到总时间中。
\end{itemize}

总体上，中位数呈现“Ping 与 TCP transfer 较小，DNS 次之，Total time 最大”的顺序。Ping 曲线集中，说明基础网络往返时间相对稳定；DNS 和 TCP transfer 更容易受服务状态影响；Total time 综合了最多因素，因此网站之间的差异最大。

\subsection{结果分析}

图中四条曲线都没有集中在 1 倍附近，说明真实互联网不可能达到只由直线距离决定的真空光速下界，主要原因包括：

\begin{enumerate}
  \item 光在光纤中的传播速度低于真空光速，而且实际光纤线路通常不是两地之间的直线。
  \item 数据包需要经过多个路由器和交换设备，每一跳都会产生处理和排队时间。
  \item DNS 查询可能需要访问缓存或上游 DNS 服务器，因此它不仅包含简单的链路传播时间。
  \item TCP 数据传输会受到网络带宽、网页大小、拥塞控制和丢包重传等因素影响。
  \item Total time 还包含 DNS、TCP 建连、TLS 握手、服务器处理、重定向和页面下载，所以其倍数最大。
\end{enumerate}

当前网站大多使用 CDN，将内容放在距离用户较近的节点上，这通常能够缩短访问时间。不过，CDN、Anycast 和 GeoIP 数据库可能使估算出的“服务器位置”与真正处理请求的位置不完全一致。因此，本实验中的 c-latency 适合用来观察总体趋势，不应当看作精确的物理测量值。

\section{实验结论}

通过本次实验，我观察到了网页访问过程中多个网络协议之间的配合。ARP 负责链路内的地址解析，DNS 将域名转换为 IP 地址，TCP 提供可靠连接，TLS 在 TCP 之上保护 HTTP 数据。Wireshark 抓包把这些原本不可见的通信步骤按数据包展示了出来。

性能测量结果表明，真实网络时延明显高于由地理距离计算出的理想下界。Ping 相对接近传播时延，而完整网页访问还包含域名解析、连接建立、加密握手、服务器处理和数据传输，因此 Total time 与 c-latency 的差距最大。实验使我对互联网的分层结构以及网络时延的来源有了较为直观的认识。

\section*{附：提交文件说明}

实验代码和测量获得的数据集附在压缩包中提交。

\end{document}
